\fancyhead[L]{\small Proof of an Algebraic Theorem}
\SourceLeaf{074}
\begin{center}
\textsc{Proof}\par\medskip
{\large\bfseries OF AN ALGEBRAIC THEOREM.}\par\medskip
\rule{0.28\linewidth}{0.4pt}
\end{center}
The subject of this paper is the theorem of \textsc{Cartes}, usually named after \textsc{Harriot}, concerning the relation between the number of positive and negative roots of an algebraic equation and the number of alternations and continuations in the signs of its coefficients. The proofs of this theorem attempted by various writers lack the clarity, brevity and comprehensive generality that one may rightly demand of so elementary a subject. A new treatment of it therefore seems not superfluous.

Let $X$ be a polynomial in $x$ of degree $m$, arranged in descending powers of $x$. We assume, without loss of generality, that its leading term is $x^m$ and that the constant term, which is independent of $x$, is not missing. Only the terms actually present are to be written down; any missing terms are therefore not to be inserted with coefficient $0$.

If the coefficients are not all positive, they will exhibit one or more changes of sign. Let $-Nx^n$ be the first negative term, $+Px^p$ the first positive term after it, $-Qx^q$ the first negative term after that, and so on. Thus $m,n,p,q$ etc. are decreasing integers, and $N,P,Q$ etc. are positive, and $X$ is represented as follows:
\[
X=x^m+\,+\, .\, .\,-Nx^n-\,-\, .\, .\,+Px^p+\,+\, .\, .\,-Qx^q-\text{etc.}
\]
\SourceLeaf{075}
Multiply $X$ by the linear factor $x-\alpha$, where $\alpha$ is assumed positive. It is easy to see that in the product $x^{n+1}$ has a negative coefficient, $x^{p+1}$ a positive one, $x^{q+1}$ a negative one, and so on, so that the product has the form
\[
X(x-\alpha)=x^{m+1}\ldots-N'x^{n+1}\ldots
+P'x^{p+1}\, .\, .\,-Q'x^{q+1}\ldots
\]
where $N',P',Q'$ etc. are positive. The signs between the displayed terms remain undetermined. But it is clear that from the leading term to the power $x^{n+1}$ there is at least one change of sign, to $x^{p+1}$ at least two, to $x^{q+1}$ at least three, and so on. If the last change of sign in $X$ occurs at the term $\pm Ux^u$, and the coefficient of $x^{u+1}$ in $X(x-\alpha)$ is denoted by $\pm U'$, then $U'$ is positive, and by the term $\pm U'x^{u+1}$ at least as many changes of sign have occurred as there are in $X$. But the last term of $X(x-\alpha)$ has the sign $\mp$, so at least one further change of sign must have occurred by then. We conclude that $X(x-\alpha)$ has at least one more change of sign than $X$.

Now let $X$ be the product of all the linear factors corresponding to the negative and imaginary roots of an equation $y=0$. Thus, if $\alpha,\beta,\gamma$ etc. are the positive roots of the same equation,
\[
y=X(x-\alpha)(x-\beta)(x-\gamma)\ldots
\]
By the preceding proposition, $X(x-\alpha)$ has at least one more change of sign than $X$, $X(x-\alpha)(x-\beta)$ at least two more, $X(x-\alpha)(x-\beta)(x-\gamma)$ at least three more, and so on. Hence, even if $X$ has no change of sign at all, $y$ has at least as many changes of sign as there are positive roots. It is evident that when the equation has neither negative nor imaginary roots we must put $X=1$, and this conclusion remains valid.

Let $y$ become $y'$ when the signs of the coefficients of the powers $x^{m-1}$, $x^{m-3}$, $x^{m-5}$ etc. are reversed. All the roots of the equation $y'=0$ are then the opposites of the roots of $y=0$. Hence $y'$ has at least as many changes of sign as the equation $y=0$ has negative roots.
\par
\SourceLeaf{076}
We therefore have the following theorem:

\emph{The equation $y=0$ cannot have more positive roots than there are changes of sign in $y$, or more negative roots than there are changes of sign in $y'$.}

This formulation of the theorem seems the most suitable, since it combines the greatest simplicity with the most comprehensive generality, and all the forms of the theorem that hold only under special conditions follow from it of themselves.

To read the bound on the number of negative roots directly from the signs of the coefficients of $y$, we must distinguish \emph{direct} sign changes and sign continuations, between terms whose exponents of $x$ differ by one, from those \emph{interrupted} by missing terms. Evidently, every direct sign change in $y'$, and every sign change interrupted by an even number of missing terms, corresponds to a sign continuation of the same kind in $y$, whereas a sign change in $y'$ interrupted by an odd number of missing terms remains a sign change of the same kind in $y$. The second part of the theorem can therefore also be expressed as follows:

The number of negative roots of the equation $y=0$ cannot exceed the number of direct sign continuations and those interrupted by an even number of missing terms, added to the number of sign changes in $y$ interrupted by an odd number of missing terms.

If no term is missing from $y$, the number of negative roots does not exceed the number of sign continuations.

If $A$ denotes the number of direct sign changes, and $B$ the number of direct sign continuations in $y$, then, when no term is missing, $A+B=m$, the number of all the roots. Since these signs tell us only that the number of positive roots cannot exceed $A$, and the number of negative roots cannot exceed $B$, it remains undecided whether there are imaginary roots, or how many. But if we know from another source that the equation has no imaginary. roots, then $A$ must necessarily equal the number of positive roots, and $B$ the number of negative roots.

\Needspace{18\baselineskip}
The situation is different when terms are missing from $y$. To see clearly
\SourceLeaf{077}
what can be concluded from this about the imaginary roots, let $a$ denote the number of sign changes interrupted by an even number of missing terms, and $c$ the number interrupted by an odd number; let $b$ and $d$ denote, respectively, the numbers of sign continuations in $y$ interrupted by an even and an odd number of missing terms. It is easy to see that $m-A-B-a-b-c-d$ equals the total number of missing terms, which we shall denote by $e$. Now, by our theorem, the number of positive roots is at most $A+a+c$, and the number of negative roots at most $B+b+c$, so the number of all the real roots is at most
\[
A+B+a+b+2c=m+c-d-e
\]
Hence the number of imaginary roots must be at least $e-c+d$.

Thus, count all the missing terms, except that in each gap with an odd number of missing terms, count one fewer at a sign change and one more at a sign continuation. The number obtained cannot exceed the number of imaginary roots.
\begin{center}\rule{0.28\linewidth}{0.4pt}\end{center}
